\section{Wstęp}
\label{sec:wstep}

\subsection{Gra Breakthrough}

\textbf{Breakthrough} to dwuosobowa gra strategiczna z pełną informacją,
wynaleziona przez Dana Troyka w 2000 roku (zwycięzca konkursu 8$\times$8 Game
Design Competition). Plansza jest kwadratowa, taka jak w szachach; każdy gracz
dysponuje pionkami ustawionymi w dwóch pierwszych rzędach po swojej stronie.
Ruch pionka odbywa się o jedno pole do przodu -- prosto lub po skosie. Ruch
prosty jest możliwy wyłącznie na puste pole; ruch po skosie może być wykonany
na pole puste lub zajęte przez pionek przeciwnika (bicie). Bicie prosto jest
niedozwolone, a pionki nie mogą się cofać.

Gra kończy się zwycięstwem gracza, który:
\begin{enumerate}[nosep]
    \item doprowadzi pionek do ostatniego rzędu przeciwnika,
    \item zbije wszystkie pionki przeciwnika,
    \item uniemożliwi przeciwnikowi wykonanie jakiegokolwiek legalnego ruchu.
\end{enumerate}
Remis nie jest możliwy.

\subsection{Dlaczego gra Breakthrough 8$\times$8}

Gra spełnia wszystkie wymagania projektu:
\begin{itemize}[nosep]
    \item \textbf{Dwuosobowa, z pełną informacją, deterministyczna} -- kluczowe
          właściwości dla metod MCTS/UCT.
    \item \textbf{Umiarkowany branching factor} -- w typowej pozycji na planszy
          8$\times$8 liczba legalnych ruchów wynosi 15--30. Głębokość partii to
          40--80 ruchów.
    \item \textbf{Proste zasady, ale trudna strategicznie} -- mimo prostych
          zasad wymaga planowania na wiele ruchów naprzód (tworzenie łańcuchów
          pionków, kontrola kolumn, poświęcenia pionka w celu wygranej).
    \item \textbf{Niszowa} -- w odróżnieniu od szachów czy warcabów gra nie jest
          aż tak popularna.
    \item \textbf{Wariant 8$\times$8 nie został rozwiązany} -- Saffidine i in.\
          \cite{saffidine2011solving} rozwiązali warianty do planszy 6$\times$5,
          a Björnsson i in.\ \cite{bjornsson2026solving} rozwiązali planszę
          6$\times$6 (w obu przypadkach wygrywa pierwszy gracz). Klasyczna
          plansza 8$\times$8 pozostaje otwarta, co zostawia miejsce na
          eksperymenty porównawcze bez znanego ,,perfekcyjnego gracza''.
\end{itemize}

\subsection{Cel projektu}

Celem projektu jest zaprojektowanie i implementacja agentów grających
w Breakthrough 8$\times$8 z wykorzystaniem algorytmu Monte Carlo Tree Search
w wariancie UCT oraz jego rozszerzeń, a następnie eksperymentalne porównanie
ich skuteczności.

W szczególności projekt obejmuje:
\begin{itemize}[nosep]
    \item implementację bazowego algorytmu UCT,
    \item implementację dwóch rozszerzeń UCT:
          \begin{itemize}[nosep]
              \item RAVE (z literatury),
              \item Progressive Bias (z literatury),
          \end{itemize}
    \item implementację gracza heurystycznego (minimax alfa-beta),
    \item porównanie jakości gry powyższych podejść,
    \item stworzenie interfejsu człowiek--komputer i przeprowadzenie testów
          z użytkownikami.
\end{itemize}

\subsection{Hipotezy badawcze}

W projekcie weryfikujemy sześć hipotez:

\begin{itemize}
    \item[\textbf{H1}] \textbf{RAVE poprawia skuteczność względem bazowego UCT.}
          Agent UCT rozszerzony o RAVE osiąga istotnie wyższy win-rate niż czysty
          UCT przy tym samym budżecie iteracji.

    \item[\textbf{H2}] \textbf{Progressive Bias poprawia skuteczność względem
          bazowego UCT.} Agent UCT z Progressive Bias osiąga istotnie wyższy
          win-rate niż czysty UCT, szczególnie przy niskich budżetach iteracji.

    \item[\textbf{H3}] \textbf{Heurystyka specjalistyczna przewyższa MCTS przy
          małym budżecie iteracji, jednak przegrywa przy wysokim.} Agent minimax
          alfa-beta z heurystyką domenową wygrywa z UCT przy małych budżetach
          iteracji, ale przegrywa przy dużych.

    \item[\textbf{H4}] \textbf{Przewaga pierwszego gracza.} W Breakthrough
          pierwszy gracz ma niewielką przewagę -- podobna sytuacja jak
          w szachach.

    \item[\textbf{H5}] \textbf{Stała eksploracji $c$ w UCT ma nieliniowy wpływ
          na siłę gry.} Istnieje optymalna wartość $c^\star$ maksymalizująca
          win-rate agenta UCT w Breakthrough 8$\times$8. Zarówno zbyt niska
          (nadmierna eksploatacja), jak i zbyt wysoka (nadmierna eksploracja)
          wartość $c$ prowadzi do istotnego spadku skuteczności.

    \item[\textbf{H6}] \textbf{Czas podejmowania decyzji MCTS rośnie w środkowej
          fazie gry.} Przy stałym budżecie iteracji średni czas na ruch agenta
          MCTS jest najwyższy w fazie środkowej partii (ruchy 15--40), gdy
          branching factor jest największy, i maleje w otwarciu oraz końcówce.
\end{itemize}

